% TOP_PROBLEM: {"id": 30006946, "title": "Uniform epsilon-equilibrium in finite multiplayer stochastic games", "category_id": 20, "category": {"id": 20, "name": "miscellaneous", "display_name": "Miscellaneous", "description": "Problems whose source classification does not fit the main mathematical categories.", "slug": "miscellaneous", "order_index": 20, "created_at": "2026-07-31T15:26:25.671Z"}, "status": "open"}
% FIRST_POSTED: 2026-09-25
% ENTRY_KIND: statement_only
% !TeX program = lualatex
% Definition and sources only; this file is not an LLM solution attempt.
\documentclass[11pt]{article}
\usepackage[margin=25mm]{geometry}
\usepackage{fontspec}
\setmainfont{Latin Modern Roman}
\usepackage{amsmath,amssymb,mathtools}
\usepackage{unicode-math}
\setmathfont{Latin Modern Math}
\usepackage{xurl,hyperref}
\hypersetup{colorlinks=true,urlcolor=blue}
\setcounter{secnumdepth}{0}
\setlength{\parindent}{0pt}
\setlength{\parskip}{5pt}
\setlength{\emergencystretch}{3em}
\begin{document}
\section*{\texorpdfstring{Uniform epsilon-equilibrium in finite multiplayer stochastic games}{Uniform epsilon-equilibrium in finite multiplayer stochastic games}}\label{30006946}
\subsection{Definitions and mathematical statement}
A finite stochastic game has finitely many players, states, and available actions, bounded stage payoffs $r_i(s,a)$, and transition probabilities $q(s'\mid s,a)$. A behavioral strategy assigns a probability distribution on the current player's actions to each observed finite history, with the information convention of the game. Given initial state $s_0$ and strategy profile $\sigma$, write
\[
 \gamma_i^n(s_0,\sigma)={\mathbb{E}}_{s_0,\sigma}
 \left[\frac1n\sum_{t=1}^n r_i(s_t,a_t)\right].
\]
The requested uniform equilibrium property is
\[
 \forall{\varepsilon}>0\ \forall s_0\ \exists\sigma\ \exists N\
 \forall n\ge N\ \forall i\ \forall\tau_i:\quad
 \gamma_i^n(s_0,\tau_i,\sigma_{-i})
 \le\gamma_i^n(s_0,\sigma)+{\varepsilon}.
\]
A deviation $\tau_i$ may itself depend on the horizon. Crucially, the same $\sigma$ works for every sufficiently long horizon. Separate finite-horizon equilibria whose strategies change with $n$ do not provide this conclusion.
\par\smallskip\textit{Formulation and concept references:} \href{https://mathconjectures.com/conjectures/GAME-001}{[S1]}.

\subsection{Short English statement}
Does every finite stochastic game admit strategies that are approximate equilibria simultaneously for all sufficiently long average-payoff horizons?
\subsection{Sources}
\begin{itemize}
\item[S1]Math Conjectures, GAME-001, catalog snapshot 31 July 2026.
\url{https://mathconjectures.com/conjectures/GAME-001}.
\textit{Formulation source.}
\end{itemize}
\end{document}
